% GENERATED FILE. Edit canonical lessons, then run npm run generate:books.
% canonical-source-hash: fnv1a64:f6c17d14368d32f8
% canonical-lessons: HI-C355-vibhag, HI-C355-arzi, HI-C355-adhikari, HI-C355-ades, HI-C355-sujhana

\chapter{Offices and Officials}
\label{ch:hi-offices-and-officials}
\hlchaptermodality{355}

\begin{chapteropening}
\textbf{By the end of this chapter:} \emph{I can deal with an office and its officials.}

Complete the last lesson of chapter 355: I can deal with an office and its officials.
\end{chapteropening}

\section[\texorpdfstring{\hi{विभाग} (vibhāg)}{विभाग (vibhāg)}]{\hi{विभाग} (vibhāg) — a department}

\label{lesson:HI-C355-vibhag}



\begin{quote}
Before the new one: say the Hindi for democracy, then the Hindi for to impose.
\end{quote}

\subsection*{You'll want to know: \hi{विभाग}}
\textbf{\hi{विभाग}} — \emph{vibhāg} — \textquotedblleft{}a department\textquotedblright{}.

\begin{quote}
\hi{मेरा} \hi{भाई} \hi{बिजली} \hi{विभाग} \hi{में} \hi{काम} \hi{करता} \hi{है।}
\end{quote}

\begin{quote}
\emph{merā bhāī bijlī vibhāg meṁ kām kartā hai.}
\end{quote}

\begin{quote}
My brother works in the electricity department.
\end{quote}

Masculine.

\begin{cousinweb}
\textbf{\hi{विभाग}} (\emph{vibhāg}) is a tatsama, borrowed from Sanskrit \emph{vibhāga}, \textquotedblleft{}a division, a share\textquotedblright{} (Monier-Williams).
\end{cousinweb}

\begin{grammarlens}[title={What you've built}]
One more word for giving reasons and opinions.
\end{grammarlens}

\subsection*{Your turn}
\begin{itemize}
\raggedright
  \item \emph{Say it:} \emph{vibhāg}
  \item \emph{Say it:} \emph{vibhāg}, once more
  \item \emph{Say it:} \emph{merā bhāī bijlī vibhāg meṁ kām kartā hai.} — My brother works in the electricity department.
  \item \emph{Recall:} say the Hindi for democracy, then the Hindi for to impose, then say \emph{vibhāg} again
\end{itemize}

\begin{tcolorbox}[breakable,colback=teal!4,colframe=teal!35!black,title={Before you move on}]
What does \hi{विभाग} mean? (\textquotedblleft{}A department\textquotedblright{}.) Say it once more.
\end{tcolorbox}
\section[\texorpdfstring{\hi{अर्ज़ी} (arzī)}{अर्ज़ी (arzī)}]{\hi{अर्ज़ी} (arzī) — an application, a petition}

\label{lesson:HI-C355-arzi}



\begin{quote}
Before the new one: say the Hindi for to impose, then the Hindi for a department.
\end{quote}

\subsection*{You'll want to know: \hi{अर्ज़ी}}
\textbf{\hi{अर्ज़ी}} — \emph{arzī} — \textquotedblleft{}an application, a petition\textquotedblright{}.

\begin{quote}
\hi{मैंने} \hi{नौकरी} \hi{के} \hi{लिए} \hi{अर्ज़ी} \hi{दी।}
\end{quote}

\begin{quote}
\emph{maiṁne naukrī ke liye arzī dī.}
\end{quote}

\begin{quote}
I applied for the job.
\end{quote}

Feminine.

\begin{cousinweb}
\textbf{\hi{अर्ज़ी}} (\emph{arzī}) is usually traced to Persian \emph{ʿarzī}, \textquotedblleft{}a petition\textquotedblright{}. The word is usually said to be built on Arabic \emph{ʿarḍ}, \textquotedblleft{}a presentation\textquotedblright{}.
\end{cousinweb}

\begin{grammarlens}[title={What you've built}]
One more word for giving reasons and opinions.
\end{grammarlens}

\subsection*{Your turn}
\begin{itemize}
\raggedright
  \item \emph{Say it:} \emph{arzī}
  \item \emph{Say it:} \emph{arzī}, once more
  \item \emph{Say it:} \emph{maiṁne naukrī ke liye arzī dī.} — I applied for the job.
  \item \emph{Recall:} say the Hindi for to impose, then the Hindi for a department, then say \emph{arzī} again
\end{itemize}

\begin{tcolorbox}[breakable,colback=teal!4,colframe=teal!35!black,title={Before you move on}]
What does \hi{अर्ज़ी} mean? (\textquotedblleft{}An application, a petition\textquotedblright{}.) Say it once more.
\end{tcolorbox}
\section[\texorpdfstring{\hi{अधिकारी} (adhikārī)}{अधिकारी (adhikārī)}]{\hi{अधिकारी} (adhikārī) — an official, an officer}

\label{lesson:HI-C355-adhikari}



\begin{quote}
Before the new one: say the Hindi for a department, then the Hindi for an application.
\end{quote}

\subsection*{You'll want to know: \hi{अधिकारी}}
\textbf{\hi{अधिकारी}} — \emph{adhikārī} — \textquotedblleft{}an official, an officer\textquotedblright{}.

\begin{quote}
\hi{अधिकारी} \hi{ने} \hi{मेरी} \hi{अर्ज़ी} \hi{पढ़ी।}
\end{quote}

\begin{quote}
\emph{adhikārī ne merī arzī paṛhī.}
\end{quote}

\begin{quote}
The official read my application.
\end{quote}

Masculine or feminine, by the person meant.

\begin{cousinweb}
\textbf{\hi{अधिकारी}} (\emph{adhikārī}) is a tatsama, borrowed from Sanskrit \emph{adhikārin}, \textquotedblleft{}a superintendent, a governor\textquotedblright{} (Monier-Williams).
\end{cousinweb}

\begin{grammarlens}[title={What you've built}]
One more word for giving reasons and opinions.
\end{grammarlens}

\subsection*{Your turn}
\begin{itemize}
\raggedright
  \item \emph{Say it:} \emph{adhikārī}
  \item \emph{Say it:} \emph{adhikārī}, once more
  \item \emph{Say it:} \emph{adhikārī ne merī arzī paṛhī.} — The official read my application.
  \item \emph{Recall:} say the Hindi for a department, then the Hindi for an application, then say \emph{adhikārī} again
\end{itemize}

\begin{tcolorbox}[breakable,colback=teal!4,colframe=teal!35!black,title={Before you move on}]
What does \hi{अधिकारी} mean? (\textquotedblleft{}An official, an officer\textquotedblright{}.) Say it once more.
\end{tcolorbox}
\section[\texorpdfstring{\hi{आदेश} (ādeś)}{आदेश (ādeś)}]{\hi{आदेश} (ādeś) — a command, an official order}

\label{lesson:HI-C355-ades}



\begin{quote}
Before the new one: say the Hindi for an application, then the Hindi for an official.
\end{quote}

\subsection*{You'll want to know: \hi{आदेश}}
\textbf{\hi{आदेश}} — \emph{ādeś} — \textquotedblleft{}a command, an official order\textquotedblright{}.

\begin{quote}
\hi{अधिकारी} \hi{ने} \hi{दफ़्तर} \hi{बंद} \hi{करने} \hi{का} \hi{आदेश} \hi{दिया।}
\end{quote}

\begin{quote}
\emph{adhikārī ne daftar band karne kā ādeś diyā.}
\end{quote}

\begin{quote}
The official ordered the office to be closed.
\end{quote}

Masculine.

\begin{cousinweb}
\textbf{\hi{आदेश}} (\emph{ādeś}) is a tatsama, borrowed from Sanskrit \emph{ādeśa}, \textquotedblleft{}an instruction; a declaration\textquotedblright{} (Monier-Williams).
\end{cousinweb}

\begin{grammarlens}[title={What you've built}]
One more word for giving reasons and opinions.
\end{grammarlens}

\subsection*{Your turn}
\begin{itemize}
\raggedright
  \item \emph{Say it:} \emph{ādeś}
  \item \emph{Say it:} \emph{ādeś}, once more
  \item \emph{Say it:} \emph{adhikārī ne daftar band karne kā ādeś diyā.} — The official ordered the office to be closed.
  \item \emph{Recall:} say the Hindi for an application, then the Hindi for an official, then say \emph{ādeś} again
\end{itemize}

\begin{tcolorbox}[breakable,colback=teal!4,colframe=teal!35!black,title={Before you move on}]
What does \hi{आदेश} mean? (\textquotedblleft{}A command, an official order\textquotedblright{}.) Say it once more.
\end{tcolorbox}
\section[\texorpdfstring{\hi{सुझाना} (sujhānā)}{सुझाना (sujhānā)}]{\hi{सुझाना} (sujhānā) — to suggest, to put (an idea) to someone}

\label{lesson:HI-C355-sujhana}



\begin{quote}
Before the new one: say the Hindi for an official, then the Hindi for a command.
\end{quote}

\subsection*{You'll want to know: \hi{सुझाना}}
\textbf{\hi{सुझाना}} — \emph{sujhānā} — \textquotedblleft{}to suggest, to put (an idea) to someone\textquotedblright{}.

\begin{quote}
\hi{अधिकारी} \hi{ने} \hi{हमें} \hi{एक} \hi{नया} \hi{रास्ता} \hi{सुझाया।}
\end{quote}

\begin{quote}
\emph{adhikārī ne hamẽ ek nayā rāstā sujhāyā.}
\end{quote}

\begin{quote}
The official suggested a new way to us.
\end{quote}

Past, with \hi{ने} (he or she …): \hi{उसने} \hi{एक} \hi{उपाय} \hi{सुझाया}, \emph{usne ek upāy sujhāyā}.

\begin{cousinweb}
\textbf{\hi{सुझाना}} (\emph{sujhānā}) is a tadbhava, inherited from Sanskrit \emph{śudhyati}, \textquotedblleft{}becomes clear\textquotedblright{} (CDIAL 12527). CDIAL glosses the Hindi \emph{sujhānā} \textquotedblleft{}to explain\textquotedblright{}; the sense of suggesting is not among its glosses, though it is close: making an idea clear to someone.
\end{cousinweb}

\begin{grammarlens}[title={What you've built}]
One more word for giving reasons and opinions.
\end{grammarlens}

\subsection*{Your turn}
\begin{itemize}
\raggedright
  \item \emph{Say it:} \emph{sujhānā}
  \item \emph{Say it:} \emph{sujhānā}, once more
  \item \emph{Say it:} \emph{adhikārī ne hamẽ ek nayā rāstā sujhāyā.} — The official suggested a new way to us.
  \item \emph{Recall:} say the Hindi for an official, then the Hindi for a command, then say \emph{sujhānā} again
\end{itemize}

\begin{tcolorbox}[breakable,colback=teal!4,colframe=teal!35!black,title={Before you move on}]
What does \hi{सुझाना} mean? (\textquotedblleft{}To suggest, to put (an idea) to someone\textquotedblright{}.) Say it once more.
\end{tcolorbox}
